%% ME_QuestionPool.tex
%% ─────────────────────────────────────────────────────────────────────────────
%% CONTROL QUESTIONS — Question Pool
%% Microeconomics, Jan Suchanek, Hochschule Bremen, Summer Semester 2026
%%
%% Covers: Ch. 2 Recap & Basics · Ch. 3 The Consumer · Ch. 4 Demand & Supply
%%         Ch. 5 Equilibrium & Welfare · Ch. 6 Market Disruptions
%%         Ch. 7 Monopolies: The Cournot Model · Ch. 9 Market Imperfections
%%         Ch. 10 State Intervention & Public Choice
%%         Ch. 12 Dominance & Strategic Interaction
%%         Ch. 13 Trade · Ch. 14 Protectionism
%%
%% Each chapter provides: 3 Multiple Choice · 2 True/False · 1 Open Question
%% ─────────────────────────────────────────────────────────────────────────────
\documentclass[paper=a5]{scrartcl}

\usepackage{xcolor}
\usepackage{hsbcommon}
\usepackage[hidelinks]{hyperref}
\usepackage[utf8]{inputenc}
\usepackage{setspace}
\usepackage{pgfplots}
\pgfplotsset{compat=1.18}

%% ── Load the question-bank package FIRST ────────────────────────────────────
%%    hsbquestions.sty must live in the SAME FOLDER as this .tex file.
%%    It defines \showanswerstrue / \showanswerSfalse via \newif\ifshowanswers.
%%    "Undefined control sequence" means the .sty was not found.
\usepackage{hsbquestions}

%% ── TOGGLE — uncomment EXACTLY ONE line ─────────────────────────────────────
\showanswerstrue               % ANSWER KEY — reveals answers + solutions
% \showanswersfalse                % STUDENT VERSION — hides answers (default)

\usepackage[headsepline]{scrlayer-scrpage}
\automark{section}
\pagestyle{scrheadings}
\ihead{\leftmark}
\chead{}
\ohead{Control Questions}

\setlength{\parskip}{1ex}

\title{Microeconomics}
\subtitle{Control Questions — Full Question Pool}
\author{Sonnet 4.6 / Claude AI using the Script of Jan \noun{Suchanek}}
% \date{Summer Semester \the\year}

\begin{document}
\setcounter{page}{247}
\onehalfspacing

\maketitle

% \begin{abstract}
% This document contains a question pool for the Microeconomics module, covering all taught chapters. Each chapter provides multiple-choice, true/false, and open questions. Toggle \texttt{\textbackslash showanswerstrue} in the preamble to reveal model answers; \texttt{\textbackslash showanswersfalse} hides them  (default, for distribution to students).
% \end{abstract}

% \tableofcontents


\clearpage

%% ═════════════════════════════════════════════════════════════════════════════
\section{Chapter 1 — Recap and Basics}
%% ═════════════════════════════════════════════════════════════════════════════

\begin{question}[type=mc]{Stock and Flow Quantities}
Which of the following is a \emph{flow} quantity?
\begin{choices}
  \choice        The total national debt of Germany.
  \choice        The oil reserves in a storage tank.
  \correctchoice The gross domestic product (GDP) per year.
  \choice        The number of cars registered in Bremen.
\end{choices}
\begin{solution}
  GDP per year is a flow quantity: it measures how much value is \emph{added}
  over a period of time (here: one year). National debt, oil reserves, and the
  number of registered cars are all stock quantities — they describe how much
  of something exists at a point in time.
\end{solution}
\end{question}

\begin{question}[type=mc]{Price Elasticity of Demand}
The price elasticity of demand $\eta$ for a good is measured at $\eta = -0.3$.
Which statement is correct?
\begin{choices}
  \choice        Demand is elastic: a 1\,\% price rise causes a 3\,\% quantity drop.
  \correctchoice Demand is inelastic: a 1\,\% price rise causes only a 0.3\,\%
                 quantity drop.
  \choice        Demand is unit elastic.
  \choice        The demand curve is horizontal.
\end{choices}
\begin{solution}
  $|\eta| = 0.3 < 1$, so demand is \textbf{inelastic}. The quantity demanded
  barely responds to the price change. Energy and petrol are classic examples
  of goods with inelastic demand.
\end{solution}
\end{question}

\begin{question}[type=mc]{Rule of 72}
An economy grows at 4\,\% per year. Approximately how many years does it take
to double in size, according to the Rule of 72?
\begin{choices}
  \choice        4 years.
  \choice        9 years.
  \correctchoice 18 years.
  \choice        72 years.
\end{choices}
\begin{solution}
  The Rule of 72 states $\Delta t \approx 72 / p$. At $p = 4\,\%$:
  $\Delta t \approx 72 / 4 = 18$ years.
\end{solution}
\end{question}

\begin{question}[type=tf]{Giffen Goods}
A Giffen good is a good for which demand \emph{rises} when its price rises.
\truefalse[true]
\begin{solution}
  \textbf{True.} Giffen goods are an exception to the law of demand. Because
  they are strongly inferior goods consuming a large share of the consumer's
  budget, the income effect of a price rise dominates the substitution effect,
  so consumers actually buy \emph{more}. Rice in subsistence economies is the
  classic textbook example.
\end{solution}
\end{question}

\begin{question}[type=tf]{Marginal vs.\ Average Values}
If the marginal value of producing one additional unit is \emph{below} the
current average, the average must be falling.
\truefalse[true]
\begin{solution}
  \textbf{True.} This is a purely mathematical relationship: whenever the next
  value added (marginal) is below the running average, it pulls the average
  down. Think of a class average: one very low new grade always lowers the mean.
\end{solution}
\end{question}

\begin{question}[type=open]{Discounting and the Interest Rate}
Explain in your own words why the interest rate can be interpreted as the
\emph{price of the future}. What happens to your valuation of a future payoff
when the interest rate rises?

\answerlines[5]

\begin{solution}
  The interest rate reflects the opportunity cost of holding money in the
  present rather than investing it. A payment received in the future is worth
  less in today's terms the higher the interest rate is, because money
  invested now would grow quickly — the present value $K_0 = K_t/(1+r)^t$
  falls as $r$ rises.

  Conversely, at near-zero interest rates, the future is valued almost as
  highly as the present (the discount factor $1/(1+r)^t \approx 1$). High
  interest rates therefore encourage present consumption and discourage
  long-term investment.
\end{solution}
\end{question}


%% ═════════════════════════════════════════════════════════════════════════════
\section{Chapter 2 — The Consumer}
%% ═════════════════════════════════════════════════════════════════════════════

\begin{question}[type=mc]{Gossen's Second Law}
A utility-maximising consumer has spent her entire budget. Gossen's Second Law
tells us that at the optimum:
\begin{choices}
  \choice        She buys only the cheapest good.
  \correctchoice The marginal utility per euro is equal across all goods purchased.
  \choice        The marginal utility of each good is zero.
  \choice        Total utility equals the budget divided by the price level.
\end{choices}
\begin{solution}
  Gossen's Second Law states that utility is maximised when the last euro
  spent on each good yields the same marginal utility. If this were not the
  case, the consumer could increase total utility by shifting spending from the
  lower-MU good to the higher-MU good.
\end{solution}
\end{question}

\begin{question}[type=mc]{Income and Substitution Effect}
The price of good $X$ rises. For a \emph{normal} good, the income effect and
the substitution effect:
\begin{choices}
  \choice        Both lead to \emph{more} consumption of $X$.
  \choice        Work in opposite directions; the income effect dominates.
  \correctchoice Both lead to \emph{less} consumption of $X$.
  \choice        Work in opposite directions; the substitution effect dominates.
\end{choices}
\begin{solution}
  For a normal good, both effects reinforce each other. The \textbf{substitution
  effect} makes $X$ relatively more expensive, so the consumer switches to
  alternatives. The \textbf{income effect} makes the consumer feel poorer, and
  since $X$ is a normal good, less of it is consumed. Both effects reduce
  the quantity demanded.
\end{solution}
\end{question}

\begin{question}[type=mc]{Inferior and Superior Goods (Engel Curves)}
As a consumer's income rises, she buys \emph{less} cheap instant noodles and
\emph{more} restaurant meals. This means:
\begin{choices}
  \choice        Both goods are normal goods.
  \correctchoice Instant noodles are an inferior good; restaurant meals are a
                 superior (luxury) good.
  \choice        Both goods are inferior goods.
  \choice        Instant noodles are a Giffen good.
\end{choices}
\begin{solution}
  A good is \textbf{inferior} if consumption falls as income rises. A good is
  \textbf{superior} (or luxury) if consumption rises faster than income.
  Instant noodles are inferior here; restaurant meals are superior.
  Their Engel curves slope in opposite directions.
\end{solution}
\end{question}

\begin{question}[type=tf]{Cobb-Douglas Utility: $\alpha + \beta = 1$}
In the Cobb-Douglas utility function $U = X^{\alpha} \cdot Y^{\beta}$ with
$\alpha + \beta = 1$, the budget shares spent on $X$ and $Y$ are constant
at $\alpha$ and $\beta$ respectively, regardless of income or prices.
\truefalse[true]
\begin{solution}
  \textbf{True.} This is a standard result for Cobb-Douglas preferences.
  Optimising the budget gives $p_X \cdot x = \alpha \cdot B$ and
  $p_Y \cdot y = \beta \cdot B$, so the fraction of the budget spent on each
  good depends only on the exponents $\alpha$ and $\beta$, not on prices or
  income. Engel curves are linear and pass through the origin.
\end{solution}
\end{question}

\begin{question}[type=tf]{Hicks Decomposition}
The Hicks decomposition isolates the substitution effect by asking: how
would the consumer re-allocate spending if relative prices changed but
\emph{real income} (utility) stayed the same?
\truefalse[true]
\begin{solution}
  \textbf{True.} In the Hicks (compensated) approach, we shift the new budget
  line outward until it is tangent to the \emph{old} indifference curve.
  The movement along that curve to the new tangency point is the pure
  substitution effect; the remaining change (between the two budget lines) is
  the income effect.
\end{solution}
\end{question}

\begin{question}[type=open]{Perfect Complements vs.\ Perfect Substitutes}
Compare perfect complements and perfect substitutes in terms of their
indifference-curve shapes, their marginal rates of substitution, and
what happens to the income and substitution effects when the price of
one good rises.

\answerlines[7]

\begin{solution}
  \textbf{Perfect complements} (e.g.\ left and right shoes) have
  L-shaped indifference curves. The MRS is zero except at the kink.
  A price rise of one good causes a pure \emph{income effect}: the
  consumer becomes poorer and buys less of both goods in fixed proportion.
  There is no substitution effect.

  \textbf{Perfect substitutes} (e.g.\ two identical shops selling fish and
  chips) have linear indifference curves. The MRS is constant ($\sigma = \infty$
  in the limit). A price rise of one good causes a pure \emph{substitution
  effect}: the consumer switches entirely to the cheaper substitute. There is
  no income effect (real utility is unchanged).

  The normal, in-between case (Cobb-Douglas) exhibits both effects partially.
\end{solution}
\end{question}


%% ═════════════════════════════════════════════════════════════════════════════
\section{Chapter 3 — Demand and Supply}
%% ═════════════════════════════════════════════════════════════════════════════

\begin{question}[type=mc]{The Supply Curve under Competition}
Under perfect competition in the short run, the firm's supply curve is:
\begin{choices}
  \choice        The average total cost curve.
  \correctchoice The marginal cost curve (above average variable cost).
  \choice        Horizontal, because all firms are price takers.
  \choice        Downward-sloping due to economies of scale.
\end{choices}
\begin{solution}
  A competitive firm sets output where $P = MC$. The supply curve is therefore
  the \textbf{marginal cost curve}. The firm only supplies if price covers
  variable costs, so the relevant portion is above average variable cost.
\end{solution}
\end{question}

\begin{question}[type=mc]{Elasticity of Supply — Short vs.\ Long Run}
Why is the price elasticity of supply typically \emph{lower} in the short run
than in the long run?
\begin{choices}
  \correctchoice In the short run, capacity is fixed; firms cannot quickly
                 expand production in response to a price rise.
  \choice        In the short run, demand is more elastic, which reduces supply
                 elasticity.
  \choice        More firms operate in the short run, reducing each firm's
                 ability to adjust.
  \choice        Fixed costs do not exist in the long run.
\end{choices}
\begin{solution}
  \textbf{Short-run supply is inelastic} because physical capacity cannot be
  changed quickly. Over the long run, new capacity can be built and new firms
  can enter, making supply much more responsive to price changes.
\end{solution}
\end{question}

\begin{question}[type=mc]{Edgeworth Box and the Trading Curve}
In an Edgeworth box with two goods and two traders, the \emph{contract curve}
(trading curve) is:
\begin{choices}
  \choice        The diagonal of the box.
  \correctchoice The locus of all points where both traders' indifference curves
                 are tangent — any move away harms at least one party.
  \choice        The locus of points where both traders are at the same utility level.
  \choice        The budget line shared by both traders.
\end{choices}
\begin{solution}
  The \textbf{contract curve} connects all Pareto-efficient allocations: points
  where the two indifference curves are tangent (equal MRS). Any point on the
  curve is stable because moving away would make at least one party worse off.
  Any point off the curve is inefficient: both parties could gain by trading.
\end{solution}
\end{question}

\begin{question}[type=tf]{The Pork Cycle}
The pork cycle arises because supply reacts to current prices with a time lag,
generating cyclical price and quantity dynamics.
\truefalse[true]
\begin{solution}
  \textbf{True.} High pork prices induce farmers to raise more pigs; supply
  increases only when the animals mature. Excess supply then depresses prices,
  triggering herd reductions — and so on. The same cobweb mechanism occurs in
  shipping, construction, and commodity markets.
\end{solution}
\end{question}

\begin{question}[type=tf]{Merit Order in Energy Markets}
In the energy merit order, the price of electricity is determined by the
marginal cost of the \emph{most expensive} technology needed to meet demand.
\truefalse[true]
\begin{solution}
  \textbf{True.} Technologies are stacked from cheapest to most expensive
  (renewables, coal, gas). The last — most expensive — technology dispatched
  sets the market price. All cheaper technologies therefore earn rents (producer
  surplus). When renewable supply is very large and covers all demand, the
  price can fall to near zero or below.
\end{solution}
\end{question}

\begin{question}[type=open]{The Pork Cycle}
Describe the pork cycle. Why does it arise, and what broader economic lesson
does it illustrate?

\answerlines[6]

\begin{solution}
  When pork prices are high, many farmers invest in raising more pigs. Because
  pigs take time to grow, supply only reaches the market with a lag of roughly
  1--2 years. The surplus then depresses prices, causing farmers to reduce
  herds. With less supply later, prices rise again — and the cycle repeats.

  The pork cycle illustrates that \textbf{markets with production lags can
  generate endogenous oscillations} (cobweb dynamics) even without external
  shocks. The same pattern underlies boom-bust cycles in shipping, construction,
  and commodity markets.
\end{solution}
\end{question}


%% ═════════════════════════════════════════════════════════════════════════════
\section{Chapter 4 — Equilibrium and Welfare}
%% ═════════════════════════════════════════════════════════════════════════════

\begin{question}[type=mc]{Consumer and Producer Surplus}
In a competitive equilibrium, the consumer surplus is:
\begin{choices}
  \choice        The area below the supply curve and above the equilibrium price.
  \correctchoice The area below the demand curve and above the equilibrium price.
  \choice        Equal to the producer surplus.
  \choice        The total revenue of all producers.
\end{choices}
\begin{solution}
  Consumer surplus is the difference between what buyers are \emph{willing to
  pay} (their reservation prices, read off the demand curve) and what they
  \emph{actually pay} (the equilibrium price). Graphically it is the triangle
  below DD and above $p^*$.
\end{solution}
\end{question}

\begin{question}[type=mc]{Welfare and Pareto Efficiency}
The outcome of a perfectly competitive market is Pareto optimal. This means:
\begin{choices}
  \choice        Everyone has the same income.
  \correctchoice No reallocation can make one party better off without making
                 another worse off.
  \choice        Total output is maximised regardless of how it is distributed.
  \choice        The government has intervened to correct market failures.
\end{choices}
\begin{solution}
  \textbf{Pareto efficiency} means we have exhausted all mutually beneficial
  trades. In a competitive equilibrium, all surplus-generating trades are
  executed, and the sum of consumer and producer surplus is maximised. It says
  nothing about the \emph{fairness} of the distribution.
\end{solution}
\end{question}

\begin{question}[type=mc]{Elasticity and Who Captures the Surplus}
Demand is highly \emph{inelastic} and supply is \emph{elastic}. Which party
captures the larger share of the total welfare gain?
\begin{choices}
  \correctchoice Consumers, because they are willing to pay far more than the
                 equilibrium price.
  \choice        Producers, because they have market power.
  \choice        The surplus is split equally.
  \choice        No surplus exists if supply is elastic.
\end{choices}
\begin{solution}
  When demand is inelastic, many consumers would have paid a much higher price,
  so the area below DD and above $p^*$ (consumer surplus) is large relative to
  the narrow triangle of producer surplus. Examples: energy, insulin.
\end{solution}
\end{question}

\begin{question}[type=tf]{Reservation Price}
A consumer's \emph{reservation price} for a good is the maximum price she
would be willing to pay for it; if the market price is below this, she earns
a positive consumer surplus.
\truefalse[true]
\begin{solution}
  \textbf{True.} Consumer surplus = reservation price $-$ market price. If the
  market price equals or exceeds the reservation price, the consumer does not
  buy (zero or negative surplus). The competitive equilibrium aggregates all
  such individual differences into total consumer surplus.
\end{solution}
\end{question}

\begin{question}[type=tf]{Pareto and Kaldor-Hicks}
A policy change that makes some people better off and others worse off
\emph{cannot} be justified under the Pareto criterion, but it \emph{can}
be justified under the Kaldor-Hicks criterion if the winners could in
principle compensate the losers.
\truefalse[true]
\begin{solution}
  \textbf{True.} The strict Pareto criterion requires that nobody be made worse
  off. Kaldor-Hicks relaxes this: a change is desirable if the gains to winners
  are large enough to \emph{potentially} compensate losers — even if no actual
  compensation takes place.
\end{solution}
\end{question}

\begin{question}[type=open]{Interpreting Consumer and Producer Surplus}
Using the demand-supply diagram, explain what consumer surplus and producer
surplus represent. Under what circumstances is one much larger than the other,
and what real-world examples illustrate this?

\answerlines[6]

\begin{solution}
  \textbf{Consumer surplus (CS)} is the sum of each buyer's reservation price
  minus the equilibrium price: the benefit consumers obtain beyond what they
  pay. Graphically it is the area below the demand curve and above $p^*$.

  \textbf{Producer surplus (PS)} is the amount above the minimum price sellers
  would accept (their marginal costs): the area above the supply curve and below
  $p^*$. If there are no fixed costs, it equals profit.

  \textbf{CS} is large relative to \textbf{PS} when demand is inelastic (e.g.\
  petrol, life-saving medicines): buyers have few substitutes and would pay
  much more than the market price.

  \textbf{PS} dominates when supply is inelastic (e.g.\ old masters paintings,
  prime real estate): the quantity is naturally limited, so sellers capture
  rents.
\end{solution}
\end{question}


%% ═════════════════════════════════════════════════════════════════════════════
\section{Chapter 5 — Market Disruptions}
%% ═════════════════════════════════════════════════════════════════════════════

\begin{question}[type=mc]{Ad Valorem vs.\ Per-Unit Tax}
An ad valorem tax on a good differs from a per-unit tax in that:
\begin{choices}
  \choice        It is paid by consumers, not producers.
  \correctchoice It is proportional to the value of the good, causing the supply
                 curve to \emph{rotate} rather than shift in parallel.
  \choice        It always generates more government revenue.
  \choice        It has no effect on consumer surplus.
\end{choices}
\begin{solution}
  A \textbf{per-unit tax} adds a fixed amount $T$ to each unit, shifting the
  supply curve upward by $T$ in parallel. An \textbf{ad valorem tax} levies a
  percentage of price, so the mark-up is larger at high prices, rotating the
  supply curve. Both can in principle produce the same equilibrium outcome
  for the right parameter choice.
\end{solution}
\end{question}

\begin{question}[type=mc]{Tax Burden and Elasticity}
A unit tax is imposed on a market where demand is highly inelastic and supply
is elastic. Who bears the greater share of the tax burden?
\begin{choices}
  \choice        Producers, because they choose the price.
  \correctchoice Consumers, because they cannot easily reduce their quantity.
  \choice        The burden is always split 50:50 regardless of elasticities.
  \choice        Nobody; the deadweight loss absorbs the entire tax.
\end{choices}
\begin{solution}
  The \textbf{less elastic} party bears more of the burden. With inelastic
  demand, consumers have few substitutes and cannot avoid the tax, so most of
  the price increase passes to them. This is why taxes on energy, tobacco, and
  alcohol are fiscally attractive: they generate stable revenue with limited
  market distortion.
\end{solution}
\end{question}

\begin{question}[type=mc]{Rigid Quantity Restriction}
A government caps the quantity of a good that may be sold to $x_2 < x^*$
(the competitive quantity). Compared to the unregulated equilibrium:
\begin{choices}
  \choice        The price falls and welfare increases.
  \choice        Only consumer surplus falls; producer surplus is unchanged.
  \correctchoice Both consumer and producer surplus change; total welfare falls
                 and a deadweight loss arises.
  \choice        The quantity restriction has no effect if demand is elastic.
\end{choices}
\begin{solution}
  A quota below $x^*$ raises the market price above $p^*$. Consumers who are
  excluded lose surplus. The net effect on producer surplus depends on the
  elasticities: some producers gain from the higher price, but others are
  forced out of the market. Total welfare falls because mutually beneficial
  trades are prevented — this is the deadweight loss triangle.
\end{solution}
\end{question}

\begin{question}[type=tf]{Supply Shocks Are Always Passed On in Full}
A cost shock (e.g.\ a sudden rise in raw-material prices) that shifts the
supply curve upward is always passed on entirely to consumers in the form
of a higher market price.
\truefalse[false]
\begin{solution}
  \textbf{False.} How much of a cost shock is passed on depends on the
  elasticities of supply and demand. The less elastic party bears more of
  the burden. In the symmetric case, the burden is shared equally. A cost
  shock will generally \emph{never} be passed on 100\,\% if demand has any
  negative slope — producers absorb part of it through lower margins.
\end{solution}
\end{question}

\begin{question}[type=tf]{Deadweight Loss from a Tax}
A tax always creates a deadweight loss — a portion of welfare that is
completely destroyed rather than transferred to the government.
\truefalse[true]
\begin{solution}
  \textbf{True.} A tax raises the price and reduces the quantity traded below
  the efficient level. The transactions that no longer take place would have
  generated surplus for both buyer and seller, but that potential gain is
  simply lost. Government revenue (GR) compensates partly, but the triangle
  of deadweight loss $DWL$ is not captured by anyone.
\end{solution}
\end{question}

\begin{question}[type=open]{Welfare Effects of a Supply Shock}
A per-unit cost increase shifts the supply curve upward. Using a supply-demand
diagram, describe how this affects price, quantity, consumer surplus, producer
surplus, and total welfare. Explain the role of elasticities.

\answerlines[7]

\begin{solution}
  The supply curve shifts up by the size of the cost increase. In the new
  equilibrium: price rises from $p_1$ to $p_2$ and quantity falls from $x_1$
  to $x_2$.

  \textbf{Consumer surplus} falls: buyers pay more and fewer units are traded.
  \textbf{Producer surplus} falls too: although price rose, their costs rose
  by more, and fewer units are sold.
  \textbf{Total welfare} falls: the new combined surplus is smaller than before,
  with a deadweight loss triangle.

  \textbf{Elasticities determine the split:} the less elastic party bears more
  of the burden. With inelastic demand (e.g.\ petrol), consumers bear most of
  the cost increase. With inelastic supply, producers bear more.
\end{solution}
\end{question}


%% ═════════════════════════════════════════════════════════════════════════════
\section{Chapter 6 — Monopolies: The Cournot Model}
%% ═════════════════════════════════════════════════════════════════════════════

\begin{question}[type=mc]{Cournot Point}
A profit-maximising monopolist sets output where:
\begin{choices}
  \choice        Price equals average cost.
  \choice        Price equals marginal cost.
  \correctchoice Marginal revenue equals marginal cost ($MR = MC$).
  \choice        Total revenue is maximised.
\end{choices}
\begin{solution}
  The monopolist maximises $\Pi(x) = R(x) - C(x)$ by setting
  $dR/dx = dC/dx$, i.e.\ $MR = MC$. This is the \textbf{Cournot point}.
  Price equals marginal cost only under perfect competition. Total revenue
  is maximised at unit elasticity, not at the profit maximum.
\end{solution}
\end{question}

\begin{question}[type=mc]{Monopoly Output vs.\ Competitive Output}
With a linear demand curve and constant marginal costs, a monopolist produces
approximately:
\begin{choices}
  \choice        The same quantity as perfect competition.
  \choice        Twice the competitive quantity.
  \correctchoice Half the competitive quantity.
  \choice        Three-quarters of the competitive quantity.
\end{choices}
\begin{solution}
  For a linear demand curve $P = a - bQ$ and constant $MC = c$, the
  competitive output is $Q_C = (a-c)/b$ and the monopoly output is
  $Q_M = (a-c)/(2b) = Q_C/2$. The monopoly price is roughly twice the
  competitive price.
\end{solution}
\end{question}

\begin{question}[type=mc]{First-Degree Price Discrimination}
A firm practising first-degree (perfect) price discrimination charges each
buyer their reservation price. The result is:
\begin{choices}
  \choice        No transactions take place because prices are too high.
  \choice        Consumer surplus is maximised at the expense of producer surplus.
  \correctchoice Consumer surplus is zero and all surplus is captured as
                 producer surplus; the quantity sold equals the competitive level.
  \choice        A larger deadweight loss than under standard monopoly pricing.
\end{choices}
\begin{solution}
  With perfect price discrimination, the monopolist extracts the full value from
  every buyer. Consumer surplus falls to zero, but every consumer who values
  the good above marginal cost still buys — so quantity equals the competitive
  level and there is no deadweight loss. In practice this requires perfect
  information about willingness to pay and no resale between buyers.
\end{solution}
\end{question}

\begin{question}[type=tf]{Monopoly and Welfare}
Compared to perfect competition, a monopoly results in a higher price, a lower
quantity traded, and a net loss of social welfare.
\truefalse[true]
\begin{solution}
  \textbf{True.} The monopolist restricts output to $Q_M < Q_C$ and charges
  $P_M > P_C$. Consumer surplus falls sharply, and the gap between the
  competitive and monopoly outcomes that is captured by neither party is the
  \textbf{deadweight welfare loss}. Producer surplus (profit) rises, but total
  welfare is lower than under competition.
\end{solution}
\end{question}

\begin{question}[type=tf]{Herfindahl-Hirschman Index (HHI)}
The Herfindahl-Hirschman Index $H = \sum_{i=1}^{N} a_i^2$ (where $a_i$ is
each firm's market share) reaches its maximum value of 1 when the market
is a perfect monopoly.
\truefalse[true]
\begin{solution}
  \textbf{True.} If there is one firm with $a_1 = 1$ (100\,\%), then
  $H = 1^2 = 1$. In perfect competition with infinitely many tiny firms, $H
  \to 0$. Competition authorities use the HHI to assess market concentration
  when evaluating mergers.
\end{solution}
\end{question}

\begin{question}[type=open]{Arguments For and Against Monopolies}
Give one economic argument \emph{in favour of} allowing a natural monopoly, and
one argument \emph{against} it. What does economic theory suggest as the best
policy response?

\answerlines[7]

\begin{solution}
  \textbf{In favour:} Natural monopolies arise where fixed costs are so high
  that a single firm can supply the entire market more cheaply than two
  competing firms (e.g.\ railway networks, electricity grids). Duplicating
  the infrastructure would waste resources.

  \textbf{Against:} An unregulated monopolist extracts rents, produces less
  than is socially optimal, and has little incentive to be efficient or
  innovate.

  \textbf{Policy:} The theoretical best response (Jean Tirole) is
  \emph{regulation} — capping the monopolist's rate of return or price —
  rather than nationalisation or break-up, since a change of ownership alone
  does not necessarily improve incentives or outcomes.
\end{solution}
\end{question}


%% ═════════════════════════════════════════════════════════════════════════════
\section{Chapter 8 — Market Imperfections}
%% ═════════════════════════════════════════════════════════════════════════════

\begin{question}[type=mc]{Adverse Selection}
Adverse selection in insurance markets arises because:
\begin{choices}
  \choice        Insurance companies invest premiums poorly.
  \correctchoice Those who most want insurance are those most likely to make a
                 claim, driving up premiums and driving low-risk customers out.
  \choice        Insured individuals start behaving more recklessly after
                 taking out a policy.
  \choice        Governments subsidise insurance companies.
\end{choices}
\begin{solution}
  \textbf{Adverse selection} (Akerlof, 1970) is a pre-contractual information
  problem: high-risk individuals have greater incentive to buy insurance than
  low-risk ones. As the insurer cannot distinguish between the two, it must
  charge a higher premium, which drives out more low-risk customers, raising
  the average risk further. In the extreme, markets can collapse entirely (the
  ``market for lemons'').
\end{solution}
\end{question}

\begin{question}[type=mc]{Pigouvian Tax}
A government introduces a Pigouvian tax equal to the marginal \emph{external}
cost of a polluting industry. The correct new equilibrium achieves:
\begin{choices}
  \choice        Zero production, since any pollution is unacceptable.
  \choice        Maximum government revenue.
  \correctchoice Production at the socially optimal level where price reflects
                 full social costs.
  \choice        An outcome identical to the untaxed market.
\end{choices}
\begin{solution}
  The Pigouvian tax internalises the externality by raising private marginal
  cost to match social marginal cost. Firms reduce output to the social
  optimum, where price reflects the full cost — including environmental damage.
  The government also collects revenue that can compensate those harmed.
\end{solution}
\end{question}

\begin{question}[type=mc]{Tragedy of the Commons}
The tragedy of the commons occurs when:
\begin{choices}
  \choice        A good is excludable but non-rival, so firms under-produce it.
  \correctchoice A resource is non-excludable and rival, so individuals over-use
                 it because private gain exceeds private cost, ignoring the
                 communal cost.
  \choice        A monopolist restricts output of a public good.
  \choice        International trade destroys a domestic commons.
\end{choices}
\begin{solution}
  Common-pool resources are \textbf{rival} (my catch reduces yours) but
  \textbf{non-excludable} (you cannot stop me fishing). Each individual's
  private benefit from over-fishing exceeds their private cost, but the
  communal cost is ignored. Solutions include privatisation, regulation, or
  — as shown by Elinor Ostrom — community governance without either.
\end{solution}
\end{question}

\begin{question}[type=tf]{Moral Hazard}
Moral hazard is an example of a \emph{post-contractual} information problem:
once insured, a person may take more risks than they would without insurance.
\truefalse[true]
\begin{solution}
  \textbf{True.} Moral hazard arises \emph{after} a contract is signed: the
  insured party's behaviour changes because they no longer bear the full
  consequences of their actions. By contrast, adverse selection is a
  \emph{pre-contractual} problem (hidden type, not hidden action).
\end{solution}
\end{question}

\begin{question}[type=tf]{Coase Theorem}
The Coase Theorem states that, provided transaction costs are zero and property
rights are clearly defined, parties will bargain to an efficient outcome
regardless of who is initially assigned the property right.
\truefalse[true]
\begin{solution}
  \textbf{True.} The key insight is that the \emph{initial} assignment of
  property rights affects distribution but not efficiency (if bargaining is
  costless). In practice, transaction costs are never zero, property rights may
  be contested, and bargaining power is often unequal — which is why Elinor
  Ostrom's work on commons governance remains important.
\end{solution}
\end{question}

\begin{question}[type=open]{Asymmetric Information and the Market for Lemons}
Explain Akerlof's ``market for lemons'' argument. Why can asymmetric
information cause a market to function poorly or even collapse? Give a
real-world example beyond used cars.

\answerlines[7]

\begin{solution}
  Akerlof observed that in the used-car market, sellers know the true quality
  of their car but buyers do not. Buyers therefore offer only the \emph{average}
  price for a random car. Owners of high-quality cars (\emph{``peaches''})
  find this average price too low and withdraw from the market. The average
  quality of cars on offer falls, so buyers lower their willingness to pay
  further — a downward spiral that can empty the market of good cars entirely.

  \textbf{Mechanisms:} Adverse selection (bad drives out good) and the absence
  of credible quality signals. Solutions include warranties, certification,
  regulation, or repeated interaction.

  \textbf{Real-world example:} Labour markets (workers know their own
  productivity better than employers), insurance (only the high-risk seek
  coverage), or credit markets (only risky borrowers may accept high interest
  rates).
\end{solution}
\end{question}


%% ═════════════════════════════════════════════════════════════════════════════
\section{Chapter 9 — The State}
%% ═════════════════════════════════════════════════════════════════════════════

\begin{question}[type=mc]{Price Ceiling}
A government sets a \emph{price ceiling below} the competitive equilibrium price.
The most likely result is:
\begin{choices}
  \choice        A surplus of goods as producers over-supply at the low price.
  \correctchoice Excess demand — more buyers want the good at the low price than
                 sellers are willing to supply.
  \choice        No change, because markets always clear at the equilibrium price.
  \choice        An increase in total welfare.
\end{choices}
\begin{solution}
  A binding price ceiling is below $p^*$. At the lower price, demand exceeds
  supply (\textbf{excess demand}). Some buyers will not find the good, creating
  queues and potentially black markets. Historical examples include Emperor
  Diocletian's bread price ceiling, which contributed to shortages.
\end{solution}
\end{question}

\begin{question}[type=mc]{Laffer Curve}
The Laffer curve illustrates that:
\begin{choices}
  \choice        Tax revenue rises monotonically with the tax rate.
  \choice        Tax revenue falls monotonically with the tax rate.
  \correctchoice Tax revenue is zero at both a 0\,\% and 100\,\% tax rate,
                 implying a revenue-maximising rate somewhere in between.
  \choice        A flat tax always maximises government revenue.
\end{choices}
\begin{solution}
  At 0\,\% tax, revenue is zero. At 100\,\%, all economic activity ceases and
  revenue is again zero. Tax revenue must therefore reach a maximum at some
  intermediate rate. The exact location of that maximum is empirically disputed
  and subject to income, substitution, and avoidance effects.
\end{solution}
\end{question}

\begin{question}[type=mc]{Ramsey Pricing Rule}
The Ramsey rule for optimal taxation states that, to minimise welfare loss,
the government should:
\begin{choices}
  \choice        Tax all goods at the same percentage rate.
  \correctchoice Tax goods with lower price elasticity more heavily, since
                 buyers cannot avoid the tax easily.
  \choice        Tax goods with higher price elasticity more heavily, to
                 maximise revenue.
  \choice        Only tax goods that are luxury items.
\end{choices}
\begin{solution}
  The Ramsey rule says to tax more heavily where the distortion is smallest,
  i.e.\ where demand (or supply) is least elastic. Highly inelastic goods
  (energy, tobacco) generate revenue with little reduction in traded quantity
  — minimising deadweight loss. A lump-sum (per-head) tax is the extreme case:
  it cannot be avoided at all, so it causes no market distortion.
\end{solution}
\end{question}

\begin{question}[type=tf]{Who Pays a Tax — Producer or Consumer?}
Whether a unit tax is legally collected from the producer or the consumer
makes no difference to the ultimate economic incidence of the tax.
\truefalse[true]
\begin{solution}
  \textbf{True.} The \emph{statutory} incidence (who writes the cheque) differs
  from the \emph{economic} incidence (who bears the cost). A tax on producers
  shifts the supply curve up; a tax on consumers shifts the demand curve down.
  Both lead to the same new equilibrium price and quantity and the same split
  of the burden between the two sides of the market.
\end{solution}
\end{question}

\begin{question}[type=tf]{Subsidies Always Increase Net Welfare}
Subsidising a market always increases total net social welfare.
\truefalse[false]
\begin{solution}
  \textbf{False.} A subsidy lowers the market price and increases quantity
  traded, raising welfare \emph{in that market}. However, the subsidy must be
  funded — through taxes elsewhere — which creates an excess burden in other
  markets. Producers also lose the incentive to reduce costs. The total
  subsidy paid exceeds the welfare gain in the subsidised market, so there is
  a net social cost.
\end{solution}
\end{question}

\begin{question}[type=open]{Price Ceilings: Why Do They Often Backfire?}
Analyse the welfare effects of a binding price ceiling below the market
equilibrium price. Explain what excess demand is, what happens to welfare,
and why black markets may emerge.

\answerlines[7]

\begin{solution}
  A binding price ceiling $\bar{p} < p^*$ reduces the quantity supplied (some
  sellers exit) while increasing quantity demanded. The result is
  \textbf{excess demand}: a gap $F-G$ between those who want to buy and those
  willing to sell at $\bar{p}$.

  \textbf{Welfare effects:} Consumer surplus is redistributed and partly
  destroyed — some buyers gain (those who can buy at the lower price) but many
  buyers are rationed. Producer surplus falls sharply. Total welfare is lower
  than at the competitive equilibrium; the lost triangle is a deadweight loss.

  \textbf{Black markets} arise because buyers willing to pay above $\bar{p}$
  and sellers willing to sell above $\bar{p}$ have an incentive to transact
  illegally. Enforcement is costly and may be ineffective, as shown
  historically by rent controls and wartime price controls.
\end{solution}
\end{question}


%% ═════════════════════════════════════════════════════════════════════════════
\section{Chapter 10 — Game Theory and Duopolies}
%% ═════════════════════════════════════════════════════════════════════════════

\begin{question}[type=mc]{Nash Equilibrium}
A Nash equilibrium in a game is a situation in which:
\begin{choices}
  \choice        Both players achieve their maximum possible payoff.
  \correctchoice No player can increase their own payoff by unilaterally
                 changing their strategy.
  \choice        Both players cooperate to maximise total welfare.
  \choice        The dominant strategy of each player is to defect.
\end{choices}
\begin{solution}
  A \textbf{Nash equilibrium} is a strategy profile in which each player's
  choice is a best response to the others' choices. No individual has an
  incentive to deviate unilaterally. Note: it need not maximise total welfare
  — the Prisoner's Dilemma demonstrates this clearly.
\end{solution}
\end{question}

\begin{question}[type=mc]{Prisoner's Dilemma}
In the one-shot Prisoner's Dilemma, the Nash equilibrium is:
\begin{choices}
  \choice        Both players cooperate.
  \correctchoice Both players defect, even though mutual cooperation would be
                 better for both.
  \choice        One player cooperates and the other defects.
  \choice        The players randomise equally between cooperate and defect.
\end{choices}
\begin{solution}
  \textbf{Mutual defection} is the Nash equilibrium: defecting is a dominant
  strategy (it yields a higher payoff regardless of the other player's choice).
  Yet mutual cooperation would be better for both — the dilemma. In repeated
  play, cooperation can be sustained by strategies such as \emph{tit for tat}.
\end{solution}
\end{question}

\begin{question}[type=mc]{Stackelberg Leadership}
In a Stackelberg duopoly, the \emph{follower} faces:
\begin{choices}
  \choice        The full market demand curve.
  \choice        A demand curve steeper than the leader's.
  \correctchoice A \emph{residual} demand curve — the part of market demand
                 not already served by the leader.
  \choice        No demand at all, since the leader captures the whole market.
\end{choices}
\begin{solution}
  The Stackelberg leader moves first and serves a portion of the market. The
  follower then acts as a de-facto monopolist over the \emph{residual demand}
  — what remains. The leader anticipates this and chooses a large quantity to
  leave little residual demand for the follower, securing a first-mover
  advantage. Example: Amazon vs.\ Otto.de.
\end{solution}
\end{question}

\begin{question}[type=tf]{Battle of the Sexes: Two Nash Equilibria}
The Battle of the Sexes game has exactly two Nash equilibria in pure
strategies.
\truefalse[true]
\begin{solution}
  \textbf{True.} In the Battle of the Sexes, Alice prefers (Concert, Concert)
  and Bob prefers (Football, Football), but both prefer to be together rather
  than apart. Both (Concert, Concert) and (Football, Football) are Nash
  equilibria in pure strategies — neither player wants to deviate unilaterally
  once they coordinate. There is also a mixed-strategy equilibrium.
\end{solution}
\end{question}

\begin{question}[type=tf]{Predatory Pricing}
Predatory pricing — setting a price so low that rivals cannot survive —
creates long-term harm to consumers even if it appears to help them in
the short run.
\truefalse[true]
\begin{solution}
  \textbf{True.} The predatory firm accepts short-run losses to drive out
  competitors, after which it raises prices to monopoly level, recouping
  losses at the consumers' expense. This is why predatory pricing is mostly
  illegal under competition law in the EU and the US.
\end{solution}
\end{question}

\begin{question}[type=open]{Why Does the Prisoner's Dilemma Matter for Trade Policy?}
Explain the logic of the Prisoner's Dilemma and apply it to the decision of
two countries over whether to impose trade tariffs on each other.

\answerlines[7]

\begin{solution}
  In the Prisoner's Dilemma, mutual cooperation yields the best collective
  outcome, but each player has a private incentive to defect regardless of
  the other's action — so rational self-interest leads to a collectively
  inferior outcome.

  \textbf{Applied to trade:} If both countries maintain free trade (cooperate),
  both gain from exchange. But if one country defects by imposing tariffs, it
  protects its domestic producers at the expense of its trading partner.
  Because both face this incentive, the Nash equilibrium is mutual tariffs —
  a trade war that leaves both countries worse off than free trade would.

  \textbf{Solution:} International agreements (WTO, bilateral treaties) act as
  commitment devices that make sustained cooperation stable, mimicking the
  ``tit for tat'' outcome in a repeated game.
\end{solution}
\end{question}


%% ═════════════════════════════════════════════════════════════════════════════
\section{Chapter 11 — Money and Trade}
%% ═════════════════════════════════════════════════════════════════════════════

\begin{question}[type=mc]{Absolute Advantage}
Country A can produce both fish and fruit more efficiently than Country B.
According to the principle of \emph{absolute} advantage:
\begin{choices}
  \choice        Trade is only beneficial for Country A.
  \choice        No trade will take place because Country B has nothing
                 to offer.
  \correctchoice Trade can still be mutually beneficial if each country
                 specialises in the good it produces relatively more efficiently.
  \choice        Country B should impose tariffs to protect its industries.
\end{choices}
\begin{solution}
  Absolute advantage says who produces a good \emph{cheaply in absolute terms};
  comparative advantage (Ricardo) says who produces it \emph{cheaply relative
  to other goods}. Even if Country A is better at producing everything, trade
  based on comparative advantage benefits both countries, as each specialises
  where its opportunity cost is lowest.
\end{solution}
\end{question}

\begin{question}[type=mc]{Comparative Advantage}
David Ricardo's principle of comparative advantage states that:
\begin{choices}
  \choice        Only countries with an absolute advantage should export goods.
  \choice        Specialisation is only beneficial if one country is better at
                 producing all goods.
  \correctchoice Each country should specialise in the good for which its
                 \emph{opportunity cost} is lowest, regardless of absolute
                 productivity levels.
  \choice        Free trade maximises welfare only when productivity is equal
                 between trading partners.
\end{choices}
\begin{solution}
  Ricardo showed that \textbf{comparative advantage} — not absolute advantage
  — drives the gains from trade. A country with lower opportunity cost in
  producing a good should export that good, even if the other country is
  absolutely more productive in every sector. Both countries end up consuming
  more than they could in autarky.
\end{solution}
\end{question}

\begin{question}[type=mc]{Trade and the Production Possibility Frontier}
When two countries trade based on comparative advantage:
\begin{choices}
  \choice        The combined production possibility frontier contracts.
  \choice        One country gains but the other loses.
  \correctchoice The combined PPF expands, allowing both countries to
                 consume beyond their individual production frontiers.
  \choice        Specialisation is complete for both countries.
\end{choices}
\begin{solution}
  Trade effectively \emph{pools} the two production frontiers. The combined
  frontier lies outside both individual frontiers, allowing consumption
  combinations otherwise unattainable. Typically, the less productive country
  specialises fully, while the more productive country may still produce some
  of both goods.
\end{solution}
\end{question}

\begin{question}[type=tf]{Who Gains and Who Loses from Trade?}
Opening a country to trade always makes every single citizen of that country
better off.
\truefalse[false]
\begin{solution}
  \textbf{False.} Free trade increases \emph{aggregate} welfare, but the gains
  and losses are distributed unevenly. In the fish-fruit example, island fruit
  producers and continental fishermen are displaced. Compensation is possible
  in principle (Kaldor-Hicks), but does not automatically occur. Political
  resistance to trade often reflects these distributional effects.
\end{solution}
\end{question}

\begin{question}[type=tf]{Comparative Advantage: Small Country}
When a small island begins trading with a large continent, the impact on the
continent's prices and trade structure is negligible, while the island
undergoes complete specialisation.
\truefalse[true]
\begin{solution}
  \textbf{True.} A small country is a price-taker on the world market: its
  trade volumes are too small to move continental prices. The island will
  fully specialise in its comparative-advantage good (e.g.\ fish), while the
  continent's relative price of fish shifts only marginally.
\end{solution}
\end{question}

\begin{question}[type=open]{Absolute vs.\ Comparative Advantage}
Distinguish between absolute and comparative advantage. Use the island-continent
fish-and-fruit example to explain why comparative advantage, not absolute
advantage, determines the pattern of trade and specialisation.

\answerlines[7]

\begin{solution}
  \textbf{Absolute advantage}: one party produces a good using fewer inputs
  than another (e.g.\ the continent is better at \emph{both} fish and fruit).

  \textbf{Comparative advantage}: one party produces a good at a lower
  \emph{opportunity cost} than another. Even if the continent is absolutely
  more efficient in both, the island may be relatively better at fishing than
  at growing fruit.

  In the example, the island trades fish at 1:4 (one fish for four fruit)
  internally, but can get the same fish at 1:1 on the continent. Island fishers
  gain dramatically by exporting. Island fruit farmers are out-competed and
  move into fishing. The island specialises totally in fish.

  \textbf{Key insight (Ricardo):} It is \emph{relative} (opportunity) costs,
  not absolute productivity, that drive trade. Both parties consume more than
  in autarky: a genuine ``gains from trade''.
\end{solution}
\end{question}


%% ═════════════════════════════════════════════════════════════════════════════
\section{Chapter 12 — Protectionism}
%% ═════════════════════════════════════════════════════════════════════════════

\begin{question}[type=mc]{Effect of an Import Tariff}
A government imposes a tariff on imports. Compared to free trade, the tariff:
\begin{choices}
  \choice        Increases consumer surplus and decreases producer surplus.
  \correctchoice Decreases consumer surplus, increases producer surplus and
                 government revenue, but creates a net deadweight loss.
  \choice        Raises the world price of the good.
  \choice        Has no effect on domestic producers.
\end{choices}
\begin{solution}
  A tariff raises the domestic price above the world price. Consumers pay more
  and buy less (\textbf{CS falls}). Domestic producers gain (\textbf{PS rises}).
  The government collects tariff revenue (\textbf{GR}). However, the total of
  CS + PS + GR is less than the original CS + PS under free trade: the two
  small triangles that are lost to nobody form the \textbf{deadweight loss}.
\end{solution}
\end{question}

\begin{question}[type=mc]{Quota vs.\ Tariff}
An import quota and a tariff can, in principle, produce the same equilibrium
price and quantity. A key difference is:
\begin{choices}
  \choice        A quota never reduces consumer surplus; a tariff always does.
  \correctchoice Under a quota, the government does not automatically receive
                 revenue; the benefit may accrue to importers as quota rents.
  \choice        A tariff affects domestic producers; a quota does not.
  \choice        Quotas are always more efficient than tariffs.
\end{choices}
\begin{solution}
  Both instruments restrict imports to the same level and produce the same
  market price. The difference is in \textbf{who captures the rents}: with a
  tariff, the government captures the revenue. With a quota, importers who
  receive licences keep the rents (the gap between world price and domestic
  price times the quota volume). If quotas are auctioned, the equivalence to
  tariffs is restored.
\end{solution}
\end{question}

\begin{question}[type=mc]{Free Trade and Welfare}
From the perspective of economic theory, free trade is:
\begin{choices}
  \choice        Always bad for small countries.
  \choice        Only beneficial if both trading partners are equally productive.
  \correctchoice Welfare-maximising in the aggregate, even though distributional
                 effects create losers within each country.
  \choice        Only beneficial when there are no trade deficits.
\end{choices}
\begin{solution}
  Theory shows that free trade maximises aggregate global welfare. Trade
  restrictions (tariffs, quotas, subsidies) create deadweight losses. However,
  the gains and losses are unevenly distributed, which generates political
  pressure for protectionism. Compensation mechanisms (redistribution) can
  in principle make everyone better off, but this rarely happens in practice.
\end{solution}
\end{question}

\begin{question}[type=tf]{Effect of a Tariff on a Small Country}
A small country imposing an import tariff cannot influence the world price;
it only raises the domestic price, so all welfare losses fall on domestic
consumers and the deadweight loss is borne entirely domestically.
\truefalse[true]
\begin{solution}
  \textbf{True.} A small country is a price taker: its tariff does not shift
  world supply and hence does not affect the world price. The tariff drives a
  wedge between the world price and the domestic price. Domestic consumers lose,
  domestic producers gain, the government gains revenue, but the net effect is a
  domestic deadweight loss. A large country can theoretically use an
  ``optimum tariff'' to extract rents from foreign exporters.
\end{solution}
\end{question}

\begin{question}[type=tf]{Export Subsidy as a Trade Policy Tool}
An export subsidy makes domestic products cheaper for foreign buyers,
increases domestic production, but reduces welfare in the subsidising country
because the cost of the subsidy exceeds the welfare gain.
\truefalse[true]
\begin{solution}
  \textbf{True.} An export subsidy shifts the supply curve down (or
  equivalently, the export supply curve). More is exported at a lower price for
  foreign buyers. Domestic consumers face a higher domestic price (since
  producers prefer to export). The net welfare effect in the subsidising country
  is negative: the subsidy cost is borne by taxpayers and exceeds the gain in
  producer surplus. Trading partners typically retaliate or challenge such
  subsidies under WTO rules.
\end{solution}
\end{question}

\begin{question}[type=open]{Why Do Countries Use Protectionist Policies Despite the
  Theoretical Case for Free Trade?}
Outline the main theoretical arguments for free trade. Then explain why
governments still resort to tariffs, quotas, and subsidies, using the
concepts of distributional effects, political economy, and the Prisoner's
Dilemma.

\answerlines[8]

\begin{solution}
  \textbf{The case for free trade:} Based on comparative advantage, trade
  expands the combined production frontier and raises aggregate welfare in all
  participating countries. Any tariff or quota creates a domestic deadweight
  loss, and retaliatory tariffs reduce welfare further for all.

  \textbf{Why protectionism persists:}
  \begin{enumerate}
    \item \textbf{Distributional effects}: Trade creates winners (competitive
      exporters) and losers (import-competing industries). Losers are often
      geographically concentrated and politically organised; winners are
      diffuse. This asymmetry creates lobbying pressure.

    \item \textbf{Political economy}: Governments face pressure from
      concentrated interest groups. Tariffs are visible, targeted transfers to
      specific industries. Subsidies serve similar functions.

    \item \textbf{Prisoner's Dilemma}: Each country has a short-term incentive
      to defect from free trade, even though mutual free trade is the
      cooperative optimum. Without binding international agreements (WTO), the
      Nash equilibrium is mutual tariffs — a trade war.
  \end{enumerate}

  \textbf{Conclusion:} International trade agreements, by making commitments
  credible, help countries overcome the Prisoner's Dilemma and sustain
  the cooperative (free-trade) equilibrium.
\end{solution}
\end{question}

\end{document}